% Cette fiche d'activité à été développée par
% + Nowar Kazem <nowar.kazem-legac@ens-rennes.fr>
% + William Hasley <william.hasley@ens-rennes.fr>

% Le template a été largement inspiré de celui de
% + Erwan Tanguy-Legac <erwan.tanguy-legac@ens-rennes.fr>
% + Héloïse Troublé <heloise.trouble@ens-rennes.fr>

% Elle est distribuée sous la licence Creative-Commons by-nc-sa 4.0
% (Attribution, Non-Commercial, Share Alike)
% qui peut être trouvée dans
% [le site de Creative Commons](https://creativecommons.org/licenses/by-nc-sa/4.0/)

\documentclass[12pt, a4paper]{article}
\usepackage[french]{babel}
\usepackage[utf8]{inputenc}
\usepackage[top=0.9in, left=0.6in, right=0.7in, bottom=1in]{geometry}
\usepackage{amssymb}
\usepackage{color}
\usepackage{fancyhdr}
\pagestyle{fancy}
\usepackage{longtable}
\usepackage{array}
\usepackage{titling}

\renewcommand{\headrulewidth}{1pt}
\chead{} 
\lhead{Les carrés magiques - Plan de l'activité}
\rhead{William Hasley et Nowar Kazem}

\renewcommand{\footrulewidth}{1pt}
\cfoot{\thepage} 
\lfoot{2025}
\rfoot{ENS Rennes}

\fancypagestyle{plain}{
	\fancyhf{} 
	\cfoot{\thepage}
	\renewcommand{\headrulewidth}{0pt}
	\renewcommand{\footrulewidth}{1pt}}

\begin{document}
	
	\title{Les carrés magiques}
	\author{William Hasley et Nowar Kazem}
	\date{31 mars 2025}
	\maketitle
	
	\paragraph*{Niveau :} CM1/CM2 
	\paragraph*{Durée :} 45 minutes
	\paragraph*{Prérequis :} Aucun
	\paragraph*{Description de l'activité :}
	Il s'agit d'une introduction aux codes correcteurs d'erreurs. Quelqu'un range des petits carrés à deux faces (une noire et l'autre blanche) dans un carré de 5x5 de 
	sorte à ce que le nombre de carrés noirs dans chaque ligne et chaque colonne soit pair. Il demande à une autre personne de tourner un carré blanc puis demande au magicien 
	de deviner le carré tourné. Avec un simple calcul on y arrive. Le but est d'expliquer aux élèves l'idée derrière ce tour de magie.
	 
	
	\paragraph*{Objectif pédagogique :}
	Expliquer aux élèves ce que sont les codes correcteurs d'erreurs et pourquoi ils sont utiles.
	
	\begin{longtable}{c|m{3.9cm}|m{8.4cm}|m{2.4cm}}		\textbf{Durée} & \textbf{Phase} & \textbf{Activités et consignes} & \textbf{Organisation et matériel} \\ \hline\hline
		
		7' &
		
		Introduction
		
		&
		
		On se présente, puis on leur dit qu'on va faire un tour de magie. On fait un exemple au tableau (on demande aux élèves de tourner un carré).
		
		& 
		
		A l'oral au tableau (avec craies ou feutres).	

		\\ \hline

		5' &

		Temps libre

		&

		On distribue le matériel (les carrés). On explique aux élèves ce qu'il doivent faire.

		&

		Groupes de 4 (les îlots), 36 carrés par groupe.
		
		\\ \hline
		
		20'
		
		&
		
		Activité
		
		& 
		
		Nous passons dans les groupes, guidons les élèves en difficulté et donnons des extensions à ceux qui vont vite. Pour les aider on peut leur remontrer le tour de magie.

		
		& Faire en sorte qu'ils comprennent le maximum.
		
		\\ \hline
		
		13' &
		
		Remise en commun + institutionnalisation
		
		&
		
		On stoppe l'activité, on reprend l'attention et on leur explique pourqoi c'est de l'informatique. On leur distribue la trace écrite et la lit ensemble.
		
		& 
		
		À l'oral, au tableau		
	
		\\ \hline
		
	\end{longtable}
	
%\newpage
	
\subsection*{Institutionnalisation :}
	
	Voir trace écrite.
	
	\paragraph*{Étayages} (si les élèves patinent) :
	\begin{itemize}
		\item leur remontrer le tour de magie, et leur montrer qu'on compte pour trouver le bon carré. 
	\end{itemize}
	
	\paragraph*{Extensions} (si certains groupes d'élèves vont trop vite) :
	\begin{itemize}
		\item utiliser un carré de 4x4 au lieu de 5x5.
		\item maintenant on fait on sorte que la somme fasse 3.
	\end{itemize}
\end{document}

